% =====================================================================
%  report.tex -- Master IEEE Conference Report
%  Project 6: Direct-Mapped Cache Controller in Verilog
%  Team: RISC Takers
%
%  Assembled and authored by Soumya Tiwari (Analysis, Synthesis &
%  Documentation Lead)
%
%  Compile on Overleaf (pdfLaTeX):
%    Upload report.tex + figures/block_diagram.pdf + figures/fsm_diagram.png
%    IEEEtran.cls is pre-installed on Overleaf — no extra upload needed.
% =====================================================================

\documentclass[conference]{IEEEtran}
\IEEEoverridecommandlockouts

% ---- Packages --------------------------------------------------------
\usepackage[T1]{fontenc}
\usepackage[utf8]{inputenc}
\usepackage{cite}
\usepackage{amsmath,amssymb,amsfonts}
\usepackage{graphicx}
\usepackage{booktabs}
\usepackage{url}
\usepackage{hyperref}
\usepackage{array}
\usepackage{multirow}
\usepackage{balance}
\usepackage{xcolor}
\usepackage{listings}
\usepackage{caption}
\usepackage{subcaption}
\usepackage{placeins}

% ---- Listings style --------------------------------------------------
\lstset{
  language=Verilog,
  basicstyle=\ttfamily\scriptsize,
  keywordstyle=\color{blue!80!black}\bfseries,
  commentstyle=\color{green!50!black}\itshape,
  stringstyle=\color{red!70!black},
  numbers=left,
  numberstyle=\tiny\color{gray},
  stepnumber=1,
  frame=single,
  breaklines=true,
  captionpos=b,
  tabsize=2
}

% ---- Hyperref setup --------------------------------------------------
\hypersetup{
  colorlinks=true,
  linkcolor=blue!60!black,
  citecolor=green!50!black,
  urlcolor=blue!70!black
}

% ---- Layout fixes (overflow / overlap) -------------------------------
% Allow line breaks after underscores so \texttt{...} names don't run
% into the neighbouring column.
\renewcommand{\_}{\textunderscore\allowbreak{}}
\pdfstringdefDisableCommands{\def\_{\textunderscore}}
\setlength{\emergencystretch}{2em}
% Float placement tuning (keeps figures/tables close to where they are cited)
\setcounter{topnumber}{3}
\setcounter{bottomnumber}{2}
\setcounter{totalnumber}{5}
\setcounter{dbltopnumber}{3}
\renewcommand{\topfraction}{0.95}
\renewcommand{\bottomfraction}{0.6}
\renewcommand{\textfraction}{0.05}
\renewcommand{\floatpagefraction}{0.8}
\renewcommand{\dbltopfraction}{0.95}
\renewcommand{\dblfloatpagefraction}{0.8}
% Ragged-right fixed-width column for tables
\newcolumntype{P}[1]{>{\raggedright\arraybackslash}p{#1}}

% ======================================================================
\begin{document}

% ---- Title -----------------------------------------------------------
\title{Design and Verification of a Parameterised\\
       Direct-Mapped Cache Controller in Verilog}

% ---- Authors ---------------------------------------------------------
\author{%
  \IEEEauthorblockN{Niharika Chaudhary}
  \IEEEauthorblockA{Cache Architecture \& Storage\\
                    Team RISC Takers}
  \and
  \IEEEauthorblockN{Karan Chillar}
  \IEEEauthorblockA{Controller FSM \& Tag Comparator\\
                    Team RISC Takers}
  \and
  \IEEEauthorblockN{Harshita Joshi}
  \IEEEauthorblockA{Memory Interface \& Top-Level Integration\\
                    Team RISC Takers}
  \and
  \IEEEauthorblockN{Shreshtha Gupta}
  \IEEEauthorblockA{Verification \& Memory-Access Traces\\
                    Team RISC Takers}
  \and
  \IEEEauthorblockN{Soumya Tiwari}
  \IEEEauthorblockA{Analysis, Synthesis \& Documentation Lead\\
                    Team RISC Takers}
}

\maketitle

% ======================================================================
% ABSTRACT
% ======================================================================
\begin{abstract}
This paper presents the design, implementation, and verification of a
fully parameterised direct-mapped cache controller for a 32-bit processor
system, implemented in synthesisable Verilog (IEEE\,1364-2001/2005). The
default configuration provides 1\,KB of data storage in 64 lines of
16\,bytes each, using a write-through, no-write-allocate policy that avoids
dirty bits and simplifies the control path. The design is partitioned into
five modular sub-systems: an address decoder, three storage RAMs (data, tag,
valid-bit), a tag comparator, a five-state synchronous finite state machine
(FSM) controller, a memory-interface bridge, and a behavioural main-memory
model. The FSM achieves single-cycle hit latency by exploiting combinational
reads from all storage arrays. Four structurally diverse memory-access traces
covering sequential, repeated-loop, random, and conflict-striding patterns are
used to evaluate performance. Bit-accurate agreement between the Python
reference model and the RTL behavioural model is confirmed across all 924
accesses. The measured hit rates are 75.00\%, 94.14\%, 0.00\%, and 0.00\%
for the sequential, loop, random, and conflict traces, yielding Average Memory
Access Time (AMAT) values of 5.00, 1.94, 17.00, and 17.00 cycles against a
16-cycle miss penalty. The complete RTL is hosted on GitHub and is intended
as the foundation for Phase~II extensions including write-back policy and
set-associativity.
\end{abstract}

\begin{IEEEkeywords}
cache controller, direct-mapped, write-through, Verilog, RTL design,
finite state machine, AMAT, trace-driven simulation, FPGA synthesis
\end{IEEEkeywords}

% ======================================================================
% SECTION I — INTRODUCTION
% ======================================================================
\section{Introduction}
\label{sec:intro}

Processor speeds have grown far faster than main-memory access times.
Without a memory hierarchy a modern CPU would spend most of its cycles
stalled on DRAM. A cache is a small, fast memory placed between the
processor and main memory that keeps recently used blocks close to the
processor, serving most accesses in a single cycle. This is effective
because real programs exhibit \emph{temporal locality}---a recently used
address is likely to be accessed again soon---and \emph{spatial
locality}---addresses near a recently used one are likely to be accessed
next~\cite{smith1982,hennessy2017}.

This project designs and verifies a parameterised direct-mapped cache
controller in Verilog. The default configuration is a 1\,KB cache with
16-byte blocks and 32-bit byte addresses. The design is partitioned into
storage (data, tag and valid-bit RAMs, address decoder), control (tag
comparator and controller FSM), and a main-memory interface. It is
exercised with sequential, repeated-loop, random, and conflict-heavy
address traces, and evaluated by hit ratio, miss ratio, AMAT, and
synthesis resource usage.

The remainder of this paper is organised as follows.
Section~\ref{sec:litrev} reviews the relevant literature.
Section~\ref{sec:solution} describes the proposed design for each
sub-system. Section~\ref{sec:results} presents the measured hit rates,
AMAT values, and synthesis estimates. Section~\ref{sec:contributions}
summarises individual contributions, and Section~\ref{sec:conclusion}
concludes with a Phase~II roadmap.

% ======================================================================
% SECTION II — LITERATURE REVIEW
% ======================================================================
\section{Literature Review}
\label{sec:litrev}

% ---- A ---------------------------------------------------------------
\subsection{Memory Hierarchy, Locality and Mapping Schemes}
\label{sec:litrev-a}

Smith's survey of cache memories~\cite{smith1982} set out the main design
choices still in use today: placement policy, block size, replacement
policy and write policy. It showed that small caches can achieve high hit
ratios because of program locality, and that block size is a trade-off:
larger blocks exploit spatial locality but raise the miss penalty and
reduce the number of distinct blocks that can be held simultaneously.

Placement policy determines where in the cache a memory block may reside.
In a \emph{direct-mapped} cache each block maps to exactly one line,
chosen by the index bits of the address. Hill and
Smith~\cite{hill1989} evaluated associativity via trace-driven simulation
and classified misses as compulsory, capacity, and conflict. Direct
mapping suffers most from conflict misses, but has a lower hit-time and
simpler hardware than associative designs because only one tag comparison
is needed. Jouppi~\cite{jouppi1990} mitigated conflict misses by adding a
small fully-associative victim cache. Hennessy and
Patterson~\cite{hennessy2017,patterson2020} unify these trade-offs into:
\begin{equation}
  \mathrm{AMAT} = t_{\mathrm{hit}} + m \cdot t_{\mathrm{penalty}},
  \label{eq:amat-basic}
\end{equation}
where $m$ is the miss rate. These sources motivate the design choices in
this project: direct mapping for a single-cycle hit path and minimal
hardware, with conflict-heavy traces to expose its main weakness.

% ---- B ---------------------------------------------------------------
\subsection{Cache Controller Design and FSM Architectures}
\label{sec:litrev-b}

The cache controller orchestrates access sequencing, hit/miss detection,
CPU stall generation, and line refills~\cite{smith1982,hennessy2017}.
Smith observed that controller complexity is strongly determined by write
and allocation policy. A \emph{write-through / no-write-allocate}
policy eliminates dirty bits and eviction write-backs, dramatically
simplifying the FSM: on a write hit the word is updated in cache and
simultaneously forwarded to memory; on a write miss the write bypasses
cache entirely.

Hennessy and Patterson~\cite{hennessy2017,patterson2020} decompose the
controller cycle into address decode, tag compare, response, and refill
allocation. The key design decision is whether tag comparison is
combinational (enabling a single-cycle hit) or pipelined. In a
direct-mapped architecture the index can be presented to asynchronous-read
RAMs immediately, so the tag and valid bit are available in the same cycle
as the request, enabling $t_{\mathrm{hit}} = 1\,\mathrm{cycle}$. Any
additional pipeline stage on the hit path adds a mandatory stall cycle to
every memory access, degrading overall processor throughput~\cite{hill1989}.

% ---- C ---------------------------------------------------------------
\subsection{Write Policies and Main-Memory Interfaces}
\label{sec:litrev-c}

The interaction between cache and main memory is governed by two orthogonal
choices: \emph{write policy} (when memory is updated) and
\emph{write-allocation policy} (whether a line is fetched on a write
miss)~\cite{smith1982,hennessy2017}. Write-through keeps main memory
always consistent, simplifying refills and system-level coherency at the
cost of elevated write-bus traffic. Write-back defers memory updates until
eviction, reducing traffic but requiring a dirty bit per line and a
two-phase eviction-refill controller~\cite{jouppi1990,patterson2020}.

Modern DRAM uses burst transfers: an initial row-activation latency is
paid for the first word, then subsequent words in the same row arrive at
one per cycle~\cite{hennessy2017}. The parameterised model in this project
uses \texttt{MEM\_FIRST\_LAT}\,=\,10 cycles and
\texttt{MEM\_NEXT\_LAT}\,=\,1 cycle. A 4-word (16-byte) block refill
therefore requires $10 + 3 \times 1 = 13$ memory cycles, which dominates
the miss penalty.

% ---- D ---------------------------------------------------------------
\subsection{Trace-Driven Simulation and Cache Verification}
\label{sec:litrev-d}

Trace-driven simulation---replaying a recorded or synthesised address
sequence through a cache model---is the standard evaluation method.
Smith~\cite{smith1982} showed it captures the interplay of block size,
associativity, and replacement policy far more faithfully than analytical
models. Hill and Smith~\cite{hill1989} used trace-driven simulation to
compare hit rates across associativities, establishing the miss
taxonomy (compulsory, capacity, conflict) used to design the trace suite
in this project. At the RTL level, cycle-accurate simulation against a
Python \emph{golden reference model} provides functional
correctness verification~\cite{patterson2020}.

% ======================================================================
% SECTION III — PROPOSED SOLUTION
% ======================================================================
\section{Proposed Solution}
\label{sec:solution}

Fig.~\ref{fig:block} shows the top-level architecture of the implemented
cache sub-system; Fig.~\ref{fig:fsm} shows the FSM state-transition
diagram. The subsections below describe each sub-system in turn.

% ---- Top-level block diagram -----------------------------------------
\begin{figure*}[!t]
  \centering
  \includegraphics[width=0.88\textwidth]{block_diagram.png}
  \caption{Top-level block diagram of the direct-mapped cache system.
           Shading indicates module ownership:
           \textit{blue} = Niharika Chaudhary (storage),
           \textit{orange} = Karan Chillar (FSM \& comparator),
           \textit{green} = Harshita Joshi (memory interface \& integration).
           The CPU-side interface exposes only six signals; all internal
           handshakes are hidden inside the \texttt{cache\_top} wrapper.}
  \label{fig:block}
\end{figure*}

% ---- A: Niharika Chaudhary -------------------------------------------
\subsection{Address Mapping and Cache Storage}
\label{sec:sol-a}

\subsubsection{Address Breakdown}

For a cache of $C$ bytes with $B$-byte blocks and an $A$-bit address, the
field widths are:
\begin{align}
  N             &= C/B,            &
  b_\mathrm{off}&= \log_2 B,       \notag \\
  b_\mathrm{idx}&= \log_2 N,       &
  b_\mathrm{tag}&= A - b_\mathrm{idx} - b_\mathrm{off}.
  \label{eq:fields}
\end{align}
With $C = 1024$\,B, $B = 16$\,B, and $A = 32$, equation~\eqref{eq:fields}
gives $N = 64$ lines, a 4-bit offset (2-bit word select + 2-bit byte
offset), a 6-bit index, and a 22-bit tag (Table~\ref{tab:addr}).

\begin{table}[!t]
  \centering
  \caption{Address field breakdown for the default cache configuration}
  \label{tab:addr}
  \renewcommand{\arraystretch}{1.25}
  \small
  \begin{tabular}{@{}lccl@{}}
    \toprule
    \textbf{Field} & \textbf{Bits} & \textbf{Width} & \textbf{Selects} \\
    \midrule
    Tag        & {[}31:10{]}  & 22 & Block identity in memory \\
    Index      & {[}9:4{]}    & 6  & One of 64 cache lines \\
    Word       & {[}3:2{]}    & 2  & One of 4 words per line \\
    Byte offset& {[}1:0{]}    & 2  & Unused (word-aligned accesses) \\
    \bottomrule
  \end{tabular}
\end{table}

\subsubsection{Storage Modules}

The cache state is split across three structures, all indexed by the
same 6-bit index and read \emph{combinationally}:

\begin{itemize}
  \item \textbf{\texttt{data\_ram}:} $64 \times 128$-bit storage
        (1\,KB). Supports a 128-bit full-line write (\texttt{we\_line})
        for cache-line refill and a 32-bit single-word write
        (\texttt{we\_word}) for write-through hits; line write takes
        priority when both are asserted simultaneously.
  \item \textbf{\texttt{tag\_ram}:} $64 \times 22$-bit storage.
        Updated synchronously on the rising edge during line refill.
  \item \textbf{\texttt{valid\_ram}:} 64-bit register vector (not a
        RAM), enabling single-cycle synchronous reset or flush of all
        valid bits. All lines are invalid after reset, guaranteeing
        compulsory misses on first access regardless of uninitialised
        tag values.
\end{itemize}

Because all three RAMs present the tag, valid bit, and data word in the
same cycle as the index, the tag comparator can resolve a hit or miss
combinationally, enabling the single-cycle fast path described in
Section~\ref{sec:sol-b}.

% ---- B: Karan Chillar ------------------------------------------------
\subsection{Controller FSM and Tag Comparator}
\label{sec:sol-b}

\subsubsection{Tag Comparator (\texttt{tag\_comparator.v})}

The tag comparator generates the \texttt{hit} signal purely
combinationally:
\begin{equation}
  \mathrm{hit} = \mathrm{rvalid} \;\land\; (\mathrm{rtag} = \mathrm{tag}).
  \label{eq:hit}
\end{equation}
Gating by \texttt{rvalid} prevents false hits on stale tag data
after reset or line invalidation: if \texttt{rvalid}\,=\,0 then
\texttt{hit}\,=\,0 regardless of the stored tag bits.

\subsubsection{Controller FSM (\texttt{cache\_controller\_fsm.v})}

Fig.~\ref{fig:fsm} and Table~\ref{tab:fsm} detail the five-state
synchronous FSM. The state encoding uses a 3-bit register
(\texttt{state[2:0]}).

\begin{figure*}[!t]
  \centering
  \includegraphics[width=0.88\textwidth]{fsm_diagram.png}
  \caption{State transition diagram of the cache controller FSM.
           The fast-path hit arc (green) completes in a single clock
           cycle from \texttt{IDLE} without entering any multi-cycle
           stall state. Read misses traverse
           \texttt{IDLE}\,$\to$\,\texttt{COMPARE\_TAG}\,$\to$\,%
           \texttt{ALLOCATE}\,$\to$\,\texttt{UPDATE}\,$\to$\,%
           \texttt{RESPOND}. Write misses skip \texttt{UPDATE} under
           the no-write-allocate policy (dashed arc).}
  \label{fig:fsm}
\end{figure*}

\begin{table*}[!t]
  \centering
  \caption{Cache Controller FSM: State Transition and Output Control Matrix}
  \label{tab:fsm}
  \renewcommand{\arraystretch}{1.3}
  \scriptsize
  \setlength{\tabcolsep}{3pt}
  \begin{tabular}{@{}llccccccccc@{}}
    \toprule
    \textbf{Current State} &
    \textbf{Condition} &
    \textbf{Next State} &
    \texttt{cpu\_ready} &
    \texttt{cpu\_stall} &
    \texttt{tag\_we} &
    \texttt{valid\_we} &
    \texttt{dwe\_word} &
    \texttt{dwe\_line} &
    \texttt{mem\_req} &
    \texttt{mem\_wr} \\
    \midrule
    \texttt{IDLE}        & no request              & \texttt{IDLE}        & 0 & 0 & 0 & 0 & 0 & 0 & 0 & 0 \\
    \texttt{IDLE}        & req, hit, read          & \texttt{IDLE}        & 1 & 0 & 0 & 0 & 0 & 0 & 0 & 0 \\
    \texttt{IDLE}        & req, hit, write         & \texttt{IDLE}        & 1 & 0 & 0 & 0 & 1 & 0 & 1 & 1 \\
    \texttt{IDLE}        & req, miss               & \texttt{COMPARE\_TAG}& 0 & 1 & 0 & 0 & 0 & 0 & 0 & 0 \\
    \midrule
    \texttt{COMPARE\_TAG}& miss                    & \texttt{ALLOCATE}    & 0 & 1 & 0 & 0 & 0 & 0 & 1 & \textit{s\_wr} \\
    \midrule
    \texttt{ALLOCATE}    & \texttt{!mem\_ready}    & \texttt{ALLOCATE}    & 0 & 1 & 0 & 0 & 0 & 0 & 1 & \textit{s\_wr} \\
    \texttt{ALLOCATE}    & \texttt{mem\_ready}, read& \texttt{UPDATE}     & 0 & 1 & 0 & 0 & 0 & 0 & 0 & 0 \\
    \texttt{ALLOCATE}    & \texttt{mem\_ready}, write& \texttt{RESPOND}   & 0 & 1 & 0 & 0 & 0 & 0 & 0 & 0 \\
    \midrule
    \texttt{UPDATE}      & (1 cycle)               & \texttt{RESPOND}     & 0 & 1 & 1 & 1 & 0 & 1 & 0 & 0 \\
    \midrule
    \texttt{RESPOND}     & (1 cycle)               & \texttt{IDLE}        & 1 & 0 & 0 & 0 & 0 & 0 & 0 & 0 \\
    \bottomrule
    \multicolumn{11}{@{}p{0.95\textwidth}@{}}{\footnotesize%
      \textit{s\_wr} = \texttt{saved\_wr} (read/write mode latched in \texttt{IDLE}).
      \texttt{dwe\_word} = \texttt{data\_we\_word};\;
      \texttt{dwe\_line} = \texttt{data\_we\_line}.}
  \end{tabular}
\end{table*}

\begin{itemize}
  \item \textbf{IDLE (\texttt{3'b000}):} Default state. If
        \texttt{cpu\_req}\,=\,1 and \texttt{hit}\,=\,1 the combinational
        fast path asserts \texttt{cpu\_ready}\,=\,1 in the \emph{same}
        cycle; for a write hit it also asserts \texttt{data\_we\_word}
        and forwards the write to memory (\texttt{mem\_req},
        \texttt{mem\_wr}).
  \item \textbf{COMPARE\_TAG (\texttt{3'b001}):} Entered on a miss from
        \texttt{IDLE}. The CPU address, write data, and
        read/write control are latched into internal registers
        (\texttt{saved\_addr}, \texttt{saved\_wdata},
        \texttt{saved\_wr}) to protect them from bus changes during
        multi-cycle stalls.
  \item \textbf{ALLOCATE (\texttt{3'b010}):} Asserts \texttt{mem\_req}.
        For a read miss the block-aligned address
        $\{\mathrm{tag},\,\mathrm{index},\,4'\mathrm{b}0000\}$ is
        driven; for a write miss the word address and data are driven
        (\texttt{mem\_wr}\,=\,1). The FSM holds this state until
        \texttt{mem\_ready}\,=\,1.
  \item \textbf{UPDATE (\texttt{3'b011}):} Active only for read misses.
        Asserts \texttt{data\_we\_line}\,=\,1,
        \texttt{tag\_we}\,=\,1, and \texttt{valid\_we}\,=\,1 for one
        cycle to write the refilled 128-bit block, update the tag RAM,
        and mark the line valid. Write misses bypass this state entirely.
  \item \textbf{RESPOND (\texttt{3'b100}):} Asserts
        \texttt{cpu\_ready}\,=\,1 for one cycle to signal completion,
        then returns unconditionally to \texttt{IDLE}.
\end{itemize}

The cycle-by-cycle latencies are:
\begin{itemize}
  \item \textbf{Read hit:} 1 cycle (combinational fast path, no state change).
  \item \textbf{Write hit:} 1 cycle (cache update + write-through dispatch in parallel).
  \item \textbf{Read miss:} $1 + 1 + 13 + 1 + 1 = 17$ cycles total
        (\texttt{IDLE}/\texttt{COMPARE\_TAG}/\texttt{ALLOCATE}/\texttt{UPDATE}/\texttt{RESPOND});
        CPU stall lasts 16 cycles.
  \item \textbf{Write miss:} $1 + 10 + 1 = 12$--$17$ cycles
        (memory write latency only; no \texttt{UPDATE}).
\end{itemize}

% ---- C: Harshita Joshi -----------------------------------------------
\subsection{Memory Interface and Top-Level Integration}
\label{sec:sol-c}

\subsubsection{Behavioural Main-Memory Model (\texttt{main\_memory.v})}

A 64\,K-word (256\,KB) flat register array pre-initialised with
$M[k] = 4k$ (word-aligned byte addresses) models DRAM timing. A
3-state internal FSM (\texttt{IDLE\,$\to$\,FIRST\,$\to$\,BURST})
counts \texttt{MEM\_FIRST\_LAT}\,=\,10 cycles for the first word,
then \texttt{MEM\_NEXT\_LAT}\,=\,1 cycle per subsequent word. When
all four words of a block are assembled, \texttt{mem\_ready} pulses
for one cycle with the full 128-bit block on \texttt{mem\_rdata}.
Write requests wait \texttt{MEM\_FIRST\_LAT} cycles then commit the
32-bit word and pulse \texttt{mem\_ready}.

\subsubsection{Memory Interface Bridge (\texttt{mem\_interface.v})}

The bridge decouples the controller FSM from raw DRAM timing. Its
two-state (\texttt{IDLE}/\texttt{BUSY}) FSM latches the outgoing
request at the \texttt{IDLE\,$\to$\,BUSY} transition, asserts
\texttt{mm\_req} for one cycle to avoid spurious re-triggering, and
routes the returned 128-bit block combinationally to
\texttt{data\_ram.wdata\_line}. The ready signal is gated:
$\texttt{fsm\_mem\_ready} = \texttt{mm\_ready} \wedge
(\texttt{mi\_state} = \texttt{BUSY})$, ensuring the FSM advances
to \texttt{UPDATE} exactly once per miss.

\subsubsection{Top-Level Wrapper (\texttt{cache\_top.v})}

\texttt{cache\_top} instantiates and connects all eight sub-modules,
exposing only the CPU-side interface:
{\small
\begin{equation}
  \begin{split}
  &\underbrace{\{\texttt{clk},\,\texttt{rst},\,\texttt{cpu\_req},\,
  \texttt{cpu\_wr},\,\texttt{cpu\_addr},\,\texttt{cpu\_wdata}\}}_{\text{inputs}}\\
  &\qquad\longrightarrow\;
  \underbrace{\{\texttt{cpu\_rdata},\,\texttt{cpu\_ready}\}}_{\text{outputs}}.
  \end{split}
  \label{eq:iface}
\end{equation}
}

% ---- D: Shreshtha Gupta ----------------------------------------------
\subsection{Verification Methodology and Trace Design}
\label{sec:sol-d}

\subsubsection{Trace Generation}

Four trace types are generated by \texttt{scripts/gen\_traces.py},
one for each class of cache miss (Table~\ref{tab:traces}):

\begin{table}[!t]
  \centering
  \caption{Trace types, access counts, and expected hit-rate ranges}
  \label{tab:traces}
  \renewcommand{\arraystretch}{1.25}
  \small
  \begin{tabular}{@{}llcc@{}}
    \toprule
    \textbf{Trace} & \textbf{Access Pattern} &
    \textbf{Accesses} & \textbf{Expected Hit Rate} \\
    \midrule
    Sequential & Linear word scan      & 200 & $\approx 75\%$ \\
    Loop       & Repeated small loop   & 324 & $\approx 94\%$ \\
    Random     & Uniform 32-bit range  & 200 & $\approx 0\%$  \\
    Conflict   & Stride = cache size   & 200 & $\approx 0\%$  \\
    \bottomrule
  \end{tabular}
\end{table}

\begin{itemize}
  \item \textbf{Sequential:} Word-by-word linear scan. Each 16-byte
        block yields one compulsory miss then three spatial-locality
        hits: steady-state hit rate $= 3/4 = 75\%$.
  \item \textbf{Loop:} A small working set (fitting in distinct cache
        lines) is traversed repeatedly. First-pass compulsory misses
        convert to 100\% temporal-locality hits on subsequent passes.
  \item \textbf{Random:} Addresses drawn uniformly from a range much
        larger than the 1\,KB cache, stressing capacity misses.
  \item \textbf{Conflict:} Pairs of addresses separated by exactly
        $1024 = 2^{10}$ bytes share the same 6-bit index, alternately
        evicting each other (ping-pong thrashing, hit rate $\approx 0\%$).
\end{itemize}

\subsubsection{Python Reference Model and Verification Flow}

A bit-accurate Python reference model (\texttt{scripts/ref\_model.py})
mirrors all RTL cache parameters. It decomposes every trace address
into tag, index, and offset, maintains a software copy of the 64-line
state (valid bit + tag), and records the expected hit or miss.
Discrepancies between the model and the RTL signal Verilog bugs, not
model errors.

The end-to-end verification flow is: (1)~generate traces; (2)~run
the reference model, logging per-access outcomes to
\texttt{results/ref\_*.log}; (3)~run the cycle-accurate behavioural
RTL simulator (\texttt{scripts/verify\_m2.py},
\texttt{verify\_all\_traces.py}); (4)~confirm bit-accurate agreement.

% ---- E: Soumya Tiwari ------------------------------------------------
\subsection{AMAT Analysis and Synthesis}
\label{sec:sol-e}

\subsubsection{AMAT Computation}

Using equation~\eqref{eq:amat-basic} with $t_{\mathrm{hit}} = 1$ cycle
and $t_{\mathrm{penalty}} = 16$ stall cycles:
\begin{align}
  t_{\mathrm{penalty}} &=
  \underbrace{1}_{\text{\scriptsize\texttt{COMPARE\_TAG}}} +
  \underbrace{13}_{\text{\scriptsize DRAM burst}} \notag \\
  &\quad +
  \underbrace{1}_{\text{\scriptsize\texttt{UPDATE}}} +
  \underbrace{1}_{\text{\scriptsize\texttt{RESPOND}}}
  = 16\;\text{stall cycles.}
  \label{eq:penalty}
\end{align}

\subsubsection{Synthesis Resource Estimation}

Yosys logic synthesis of the M1 storage modules confirms the expected
register bit counts. Full FPGA synthesis (Xilinx Artix-7
\texttt{xc7a35t}) is targeted in Phase~II; estimated resource usage is
shown in Table~\ref{tab:synth}.

\begin{table}[!t]
  \centering
  \caption{Estimated FPGA resource utilisation (Xilinx Artix-7 target)}
  \label{tab:synth}
  \renewcommand{\arraystretch}{1.25}
  \small
  \begin{tabular}{@{}lcP{3.1cm}@{}}
    \toprule
    \textbf{Resource} & \textbf{Count / Value} & \textbf{Notes} \\
    \midrule
    Data RAM bits    & 8\,192  & $64 \times 128\,\mathrm{b}$ \\
    Tag RAM bits     & 1\,408  & $64 \times 22\,\mathrm{b}$ \\
    Valid-bit register & 64   & 64-bit vector; single-cycle flush \\
    FSM state reg.   & 3\,FF  & \texttt{3'b000}--\texttt{3'b100} \\
    LUTs (est.)      & ${\sim}80$ & Decoder + comparator + FSM comb.\ \\
    FFs  (est.)      & ${\sim}190$& RAMs + FSM + request latch \\
    $F_{\max}$ (est.)& $>100\,\mathrm{MHz}$ & Short comb.\ hit path \\
    \bottomrule
  \end{tabular}
\end{table}

% ======================================================================
% SECTION IV — RESULTS AND DISCUSSION
% ======================================================================
\section{Results and Discussion}
\label{sec:results}

\subsection{Hit Rate and AMAT Results}

Table~\ref{tab:amat} presents the hit rates, miss rates, and AMAT values
measured by the Python reference model and confirmed bit-accurately by
the cycle-accurate RTL behavioural model across all 924 accesses
(4/4 traces, 100\% agreement).

\begin{table}[!t]
  \centering
  \caption{Measured hit/miss rates and computed AMAT for all four traces
           ($t_{\mathrm{hit}} = 1$ cycle, $t_{\mathrm{penalty}} = 16$ cycles)}
  \label{tab:amat}
  \renewcommand{\arraystretch}{1.4}
  \resizebox{\columnwidth}{!}{%
  \begin{tabular}{@{}lcccccc@{}}
    \toprule
    \textbf{Trace} &
    \textbf{Accesses} &
    \textbf{Hits} &
    \textbf{Misses} &
    \textbf{Hit Rate} &
    \textbf{Miss Rate} &
    \textbf{AMAT (cycles)} \\
    \midrule
    Sequential
      & 200 & 150 &  50
      & 75.00\% & 25.00\%
      & $1 + 0.25 \times 16 = \mathbf{5.00}$ \\
    Loop
      & 324 & 305 &  19
      & 94.14\% &  5.86\%
      & $1 + 0.0586 \times 16 = \mathbf{1.94}$ \\
    Random
      & 200 &   0 & 200
      &  0.00\% & 100.00\%
      & $1 + 1.00 \times 16 = \mathbf{17.00}$ \\
    Conflict
      & 200 &   0 & 200
      &  0.00\% & 100.00\%
      & $1 + 1.00 \times 16 = \mathbf{17.00}$ \\
    \midrule
    \textbf{Total / Weighted avg.}
      & \textbf{924} & \textbf{455} & \textbf{469}
      & 49.24\% & 50.76\%
      & --- \\
    \bottomrule
  \end{tabular}}
\end{table}

\subsection{Discussion}

\textbf{Sequential trace (AMAT = 5.00\,cycles).}
A linear word-by-word scan incurs one compulsory miss per 16-byte block
followed by three spatial-locality hits ($3/4 = 75\%$ hit rate). The
AMAT of 5.00 cycles confirms that the cache effectively captures
spatial locality: performance is $17/5 = 3.4\times$ better than the
all-miss baseline.

\textbf{Loop trace (AMAT = 1.94\,cycles).}
A small working set that fits within distinct cache lines incurs
compulsory misses only on the first pass; every subsequent access is a
temporal-locality hit. The resulting 94.14\% hit rate gives an AMAT of
1.94\,cycles---near-ideal performance and only $1.94\times$ the
theoretical minimum of 1\,cycle.

\textbf{Random trace (AMAT = 17.00\,cycles).}
Random addresses drawn from the full 32-bit space yield a 0\% hit rate:
the probability that a random access finds a valid matching line in
only 64 cache lines is negligible. The AMAT equals the all-miss
worst case of 17\,cycles. This highlights that a 1\,KB direct-mapped
cache is ineffective against workloads with no temporal or spatial
locality.

\textbf{Conflict trace (AMAT = 17.00\,cycles).}
Addresses separated by exactly $2^{10}$\,bytes ($=1\,$KB) share the
same 6-bit index. Alternating between two such addresses causes each
access to evict the line loaded by the previous one (ping-pong
thrashing), driving the hit rate to exactly 0\% and AMAT to its maximum
value. This is the defining pathology of direct-mapped caches and
directly motivates 2-way set-associativity in Phase~II.

\subsection{Verification Summary}

All correctness checks pass with no failures:
\begin{itemize}
  \item \textbf{Storage unit tests (Niharika Chaudhary):}
        2\,012 address-decoder field checks; 266 data-RAM word read-backs
        (including partial writes and priority); 68 tag-RAM checks; 196
        valid-bit checks (reset, flush, write-priority).
  \item \textbf{FSM verification (Karan Chillar):}
        97/97 functional assertions covering reset, read hit, read miss,
        write hit, write miss, and back-to-back accesses.
  \item \textbf{M1+FSM pipeline integration (Karan Chillar):}
        6/6 scenarios: cold miss, spatial locality, write hit, temporal
        hit, conflict miss, and re-access after eviction.
  \item \textbf{Trace verification (Shreshtha Gupta):}
        4/4 traces, 924 accesses, 100\% bit-accurate match.
  \item \textbf{Full-system integration (Harshita Joshi):}
        \texttt{tb\_cache\_top.v} validates cold misses, spatial locality,
        write-through, no-write-allocate, and conflict-miss thrashing
        over the complete 8-module \texttt{cache\_top} hierarchy.
\end{itemize}

% ======================================================================
% SECTION V — INDIVIDUAL CONTRIBUTIONS
% ======================================================================
\section{Individual Contributions}
\label{sec:contributions}

\begin{table}[!t]
  \centering
  \caption{Team member contributions}
  \label{tab:contrib}
  \renewcommand{\arraystretch}{1.3}
  \footnotesize
  \begin{tabular}{@{}P{1.9cm}P{1.9cm}P{4.1cm}@{}}
    \toprule
    \textbf{Name} & \textbf{Area} & \textbf{Key Deliverables} \\
    \midrule
    Niharika Chaudhary &
      Storage \& Architecture &
      \texttt{address\_decoder.v}, \texttt{data\_ram.v},
      \texttt{tag\_ram.v}, \texttt{valid\_ram.v},
      \texttt{interfaces.md}, block diagram, unit testbenches \\
    \midrule
    Karan Chillar &
      Controller FSM \& Comparator &
      \texttt{tag\_comparator.v},
      \texttt{cache\_controller\_fsm.v},
      FSM diagram, Python verification suite,
      trace verification \\
    \midrule
    Harshita Joshi &
      Memory Interface \& Integration &
      \texttt{main\_memory.v},
      \texttt{mem\_interface.v},
      \texttt{cache\_top.v},
      integration testbenches \\
    \midrule
    Shreshtha Gupta &
      Verification \& Traces &
      \texttt{gen\_traces.py},
      \texttt{ref\_model.py},
      4 trace files,
      \texttt{cache\_tb.v} \\
    \midrule
    Soumya Tiwari &
      Analysis \& Documentation &
      \texttt{analysis.py},
      AMAT tables, synthesis estimate,
      master IEEE report (\texttt{report.tex}),
      README \\
    \bottomrule
  \end{tabular}
\end{table}

Table~\ref{tab:contrib} details each team member's ownership and
deliverables.

% ======================================================================
% SECTION VI — CONCLUSION
% ======================================================================
\FloatBarrier
\section{Conclusion}
\label{sec:conclusion}

This paper has presented a fully parameterised, synthesisable
direct-mapped cache controller implemented in Verilog and verified
through trace-driven simulation. All primary design goals have been
met: a single-cycle hit latency via combinational tag comparison, a
clean five-state FSM implementing write-through / no-write-allocate
policy, and bit-accurate functional correctness across a diverse
four-trace benchmark suite.

The experimental results quantitatively confirm theoretical predictions.
The loop trace achieves near-ideal AMAT of 1.94\,cycles under temporal
locality, while the conflict trace reveals the defining weakness of
direct mapping (AMAT = 17.00\,cycles under ping-pong thrashing). These
findings directly motivate the following Phase~II extensions:

\begin{enumerate}
  \item \textbf{2-way set-associativity} with LRU replacement to
        eliminate conflict misses at the cost of a simple replacement
        policy and one additional tag comparison.
  \item \textbf{Write-back policy} with per-line dirty bits to reduce
        memory-bus traffic, requiring additional eviction states in the
        FSM.
  \item \textbf{Parameter sweep} across cache sizes (512\,B to 8\,KB)
        and block sizes (8\,B to 64\,B), generating AMAT vs.\ resource
        trade-off curves using Yosys and Xilinx Vivado synthesis.
\end{enumerate}

The complete RTL, testbenches, verification scripts, and traces are
publicly available at
\url{https://github.com/NullInTheHead/risc-takers-cache-controller},
enabling full reproducibility of all results reported in this paper.

% ======================================================================
% REFERENCES
% ======================================================================
\newpage
\begin{thebibliography}{9}

\bibitem{smith1982}
A.~J. Smith, ``Cache memories,''
\emph{ACM Computing Surveys},
vol.~14, no.~3, pp.~473--530, Sep.\ 1982.

\bibitem{hill1989}
M.~D. Hill and A.~J. Smith, ``Evaluating associativity in CPU caches,''
\emph{IEEE Transactions on Computers},
vol.~38, no.~12, pp.~1612--1630, Dec.\ 1989.

\bibitem{jouppi1990}
N.~P. Jouppi, ``Improving direct-mapped cache performance by the addition
of a small fully-associative cache and prefetch buffers,'' in
\emph{Proc.\ 17th Annual International Symposium on Computer Architecture
(ISCA)}, Pittsburgh, PA, USA, 1990, pp.~364--373.

\bibitem{hennessy2017}
J.~L. Hennessy and D.~A. Patterson,
\emph{Computer Architecture: A Quantitative Approach}, 6th~ed.\
Cambridge, MA, USA: Morgan Kaufmann, 2017.

\bibitem{patterson2020}
D.~A. Patterson and J.~L. Hennessy,
\emph{Computer Organization and Design RISC-V Edition:
The Hardware Software Interface}, 2nd~ed.\
San Francisco, CA, USA: Morgan Kaufmann, 2020.

\end{thebibliography}

\balance   % balance the two columns on the last page

\end{document}